% DEPRECATED 2026-10-31: 本文件已不再被 main.tex include；内容已并入 6_vision_encoder.tex (§6) / 5_llm_posttrain.tex (§5.7) / 10_conclusion.tex (§10)。保留仅供历史对照。
\section{MLLM Alignment}\label{sec:mllm}

Stage~(iv) connects the vision encoder of Section~\ref{sec:vision} to the language backbone through a lightweight projector, following the staged curriculum established for open multimodal assistants: projector pre-training, mid-training, and instruction tuning.
% 涓枃锛氱 (iv) 闃舵閫氳繃涓€涓交閲忔姇褰卞眰锛屾妸绗?\ref{sec:vision} 鑺傜殑瑙嗚缂栫爜鍣ㄦ帴鍒拌瑷€楠ㄥ共缃戠粶涓婏紝娌跨敤涓哄紑鏀惧妯℃€佸姪鎵嬪缓绔嬬殑鍒嗛樁娈佃绋嬶細鎶曞奖灞傞璁粌銆佷腑鏈熻缁冧笌鎸囦护寰皟銆?
The design is motivated by the observation that the dominant EDA artifacts are visual---layout and schematic views account for $70.2\%$ of the tasks in the evaluation benchmark~\cite{liu2026zhulong}---so a multimodal core can be scored on the same execution-based protocol as the text-only one.
% 涓枃锛氳璁捐鐨勫姩鏈哄湪浜庯細EDA 鐨勪富瑕佷骇鐗╂槸瑙嗚鎬х殑鈥斺€斿湪璇勬祴鍩哄噯 \cite{liu2026zhulong} 涓紝鐗堝浘涓庡師鐞嗗浘瑙嗗浘鍗犱换鍔＄殑 $70.2\%$鈥斺€斿洜姝ゅ妯℃€佸唴鏍稿彲浠ヤ笌绾枃鏈唴鏍搁噰鐢ㄥ悓涓€濂楀熀浜庢墽琛岀殑鍗忚鎵撳垎銆?

\paragraph{Task construction.}
We build three task families, each chosen because it admits an objective ground truth.
% 涓枃锛氫换鍔℃瀯寤恒€傛垜浠瀯寤轰笁绫讳换鍔★紝姣忎竴绫荤殑鍏ラ€夌悊鐢遍兘鏄畠鍏峰瀹㈣鐨勬爣鍑嗙瓟妗堛€?
(i) \emph{image$\to$script}: given a rendered layout or schematic view, produce the PyAether or SKILL script that recreates it; supervision is available because executing a script yields its own reference rendering.
% 涓枃锛?i) \emph{鍥?\to$鑴氭湰}锛氱粰瀹氭覆鏌撳嚭鐨勭増鍥炬垨鍘熺悊鍥捐鍥撅紝鐢熸垚鑳介噸鐜板畠鐨?PyAether 鎴?SKILL 鑴氭湰锛涚洃鐫ｄ俊鍙峰ぉ鐒跺彲寰楋紝鍥犱负鎵ц鑴氭湰鏈韩灏变細浜х敓鍙傝€冩覆鏌撳浘銆?
(ii) \emph{script$\to$image}: predict the geometry that a given script produces, again self-labelled by execution.
% 涓枃锛?ii) \emph{鑴氭湰$\to$鍥緘锛氶娴嬬粰瀹氳剼鏈墍浜х敓鐨勫嚑浣曞浘褰紝鍚屾牱鐢辨墽琛岃嚜鏍囨敞銆?
(iii) \emph{layout VQA}: answer attribute questions---layer identity, device type, spacing and rule violations---labelled by domain experts.
% 涓枃锛?iii) \emph{鐗堝浘 VQA}锛氬洖绛斿睘鎬х被闂鈥斺€斿眰娆℃爣璇嗐€佸櫒浠剁被鍨嬨€侀棿璺濅笌瑙勫垯杩濅緥鈥斺€旂敱棰嗗煙涓撳鏍囨敞銆?
The first family is the strongest of the three: its output is a script, so it is scored by the same assertion sandbox used in stage~(ii), which gives the multimodal and text-only stages a common metric instead of two incommensurable ones.
% 涓枃锛氱涓€绫绘槸涓夎€呯殑鏈€寮洪」锛氬叾杈撳嚭鏄剼鏈紝鍥犳鍙敤涓庣 (ii) 闃舵鐩稿悓鐨勬柇瑷€娌欑洅璇勫垎锛岃繖浣垮妯℃€侀樁娈典笌绾枃鏈樁娈靛叡浜悓涓€搴﹂噺锛岃€屼笉鏄悇鐢ㄤ竴濂楁棤娉曟崲绠楃殑鎸囨爣銆?

\paragraph{Curriculum and ablations.}
We run stage-1 projector pre-training on caption-style pairs, stage-1.5 mid-training on the three families above, and stage-2 instruction tuning on the full mixture.
% 涓枃锛氳绋嬩笌娑堣瀺銆傛垜浠厛鐢ㄨ鏄庢枃瀛楀紡鐨勫浘鏂囧鍋?stage-1 鎶曞奖灞傞璁粌锛屽啀鐢ㄤ笂杩颁笁绫讳换鍔″仛 stage-1.5 涓湡璁粌锛屾渶鍚庡湪鍏ㄩ噺娣峰悎鏁版嵁涓婂仛 stage-2 鎸囦护寰皟銆?
The ablations vary the projector (MLP versus query-based), whether the vision tower is frozen or updated, and the number of visual tokens per image.
% 涓枃锛氭秷铻嶅彉閲忓寘鎷細鎶曞奖灞傜被鍨嬶紙MLP 涓庡熀浜?query 鐨勬柟妗堬級銆佽瑙夊鏄惁鍐荤粨鎴栨洿鏂般€佷互鍙婃瘡寮犲浘鍍忕殑瑙嗚 token 鏁般€?
The last is of particular interest.
% 涓枃锛氬叾涓渶鍚庝竴椤瑰挨鍊煎緱鍏虫敞銆?
The architecture study of Section~\ref{sec:vision} selected $224/16$, where an image becomes only $14{\times}14 = 196$ tokens (${\sim}16{\times}16$\,px/token)---a regime in which the training objective is dominated by coarse semantics. For detail-dense line drawings, a $336$-px view would yield $576$ tokens ($2.9{\times}$ more) and a $448$-px view $1024$ tokens ($5.2{\times}$ more); whether the $196$-token setting transfers to such views is precisely what stage~(iv) tests.
% 涓枃锛氱 \ref{sec:vision} 鑺傜殑鏋舵瀯鐮旂┒閫夊畾浜?$224/16$锛屾鏃朵竴寮犲浘鍍忎粎鍙樹负 $14{\times}14 = 196$ 涓?token锛堢害 $16{\times}16$ 鍍忕礌/token锛夆€斺€斿湪杩欎竴鍖洪棿锛岃缁冪洰鏍囪绮楃矑搴﹁涔変富瀵笺€傚浜庣粏鑺傚瘑闆嗙殑绾跨锛?336$ 鍍忕礌瑙嗗浘浼氬緱鍒?$576$ 涓?token锛堝 $2.9$ 鍊嶏級锛?448$ 鍍忕礌瑙嗗浘鍒欎负 $1024$ 涓紙澶?$5.2$ 鍊嶏級锛?196$-token 鐨勮瀹氳兘鍚﹁縼绉诲埌杩欑被瑙嗗浘锛屾鏄 (iv) 闃舵瑕佹楠岀殑銆?
Layout and schematic views are the opposite: they are detail-dense line drawings whose meaning lives in fine geometry.
% 涓枃锛氱増鍥句笌鍘熺悊鍥捐鍥炬伆鎭扮浉鍙嶏細瀹冧滑鏄粏鑺傚瘑闆嗙殑绾跨锛屽叾璇箟瀛樺湪浜庣簿缁嗙殑鍑犱綍缁撴瀯涔嬩腑銆?
We therefore treat resolution and visual-token count as first-order variables here rather than inheriting the earlier null result, which we read as a property of the proxy task and not of the encoder.
% 涓枃锛氬洜姝ゆ垜浠湪姝ゆ妸鍒嗚鲸鐜囦笌瑙嗚 token 鏁板綋浣滀竴闃跺彉閲忓鐞嗭紝鑰屼笉鐩存帴缁ф壙姝ゅ墠鐨勯浂缁撴灉鈥斺€旀垜浠皢鍏剁悊瑙ｄ负浠ｇ悊浠诲姟鐨勫睘鎬э紝鑰岄潪缂栫爜鍣ㄧ殑灞炴€с€?

Results are reported as \textbf{[TBD]}.
% 涓枃锛氱粨鏋滀互 \textbf{[TBD]} 琛ㄧず锛堝緟濉級銆
